\section{Related Work}
\label{sec:related}

The BabyLM Challenge has run three times, and its findings reports serve
as a census of what helps under developmental budgets
\citep{warstadt2023findings,hu2024findings,charpentier2025findings}. The
recipes that succeed tend to change what the model predicts, not
when. GPT-BERT, our backbone, won the 2024 edition by blending masked and
causal prediction in a single transformer \citep{charpentier2024gptbert}.
In the 2025 edition, reweighting the masking distribution as training
progresses outperformed the official strict-small baseline
\citep{edman2025mask}. BabyLlama distilled an ensemble
of teachers into one small student and later served as an official
baseline \citep{timiryasov2023babyllama,hu2024findings}. All of these run
as a single uninterrupted phase; where in the budget the signal lands is
not itself a design variable. That variable is our subject.

Timing has been studied descriptively: \citet{chang2022word} chart when
individual words are acquired during pretraining, regularly enough to
compare against child norms, and the 2026 evaluation includes an
age-of-acquisition task built on this lens \citep{choshen2026babylm}. We
push the lens from describing when learning happens to acting on when it
stops.

The interventions in our controlled study (\S\ref{sec:search}) each arrive
with a published pedigree: span corruption \citep{joshi2020spanbert},
selective token losses that concentrate compute on learnable,
worth-learning points \citep{mindermann2022prioritized,lin2024rho}, the
Muon optimizer \citep{jordan2024muon}, weight averaging in its stochastic
and soup forms \citep{izmailov2018swa,wortsman2022soups}, and
morpheme-aware segmentation from a child-inspired entry
\citep{bolucu2025morpheme}. The reported gains come from other regimes:
larger models, longer budgets, a previous edition's setting. Whether
they hold under a 10M-word budget is empirical, and \S\ref{sec:search}'s
answer is that none does. What informs our method is not the failures but
the pattern in how they fail.

Preference optimization entered the toolkit as an alignment method: DPO
tunes a model to prefer one completion over another under a KL penalty
toward a frozen reference, with pairs elicited from human judgment
\citep{rafailov2023dpo}. Its one BabyLM appearance so far targets
dialogue: \citet{padovani2025dialogue} tune a dialogue-pretrained model
with DPO toward communicative continuations, gaining on a custom dialogue
benchmark while conceding the standard suite. \MPO{} retains the objective
and replaces everything around it. Pairs are a corpus sentence and a
synthetic corruption of it, so no annotator, reward model, or teacher is
involved; the target is grammatical competence, not helpfulness; and the
phase is placed by measurement at the saturation checkpoint, not appended
to a finished model. The corruption
design itself borrows from minimal-pair evaluation, in which BLiMP
isolates one phenomenon per item by contrasting two strings that differ
minimally \citep{warstadt2020blimp}. Learning from corruptions also has a
pretraining pedigree: ELECTRA trains a token-level discriminator against
sampled replacements throughout pretraining \citep{clark2020electra},
where \MPO{} compares whole sequences, after acquisition, under an
anchor. We train on the form of comparison the benchmarks measure,
computed entirely on text the rules already allow.
